\chapter{Bloch--Kato Selmer Groups, Berger--Wach Arcs, and Twisted-Place Integrality}
\label{ch:18}

Part~VI closes the fourth and final arithmetic pillar of our Arakelov-filtered automorphic period tower (\cref{def:01:period-tower}(v)): controlling the \emph{twisted-place integrality deficit} $\mathcal{D}_{\mathrm{tw}}(x)$ that enters the lower face \eqref{eq:06:lower-face} of the Two-Sided Arakelov Slope Sandwich and the definite-side closure (\cref{thm:17:definite-assembly}).

Whenever a rational prime $\ell$ is split or twisted in the host CM extension $E/\QQ$ (or ramifies at the special primes $\ell \in \{2, 3\}$ already handled in \cref{thm:17:at23}), the definite unitary period kernels acquire local $\ell$-adic denominators arising from congruences between the attached $\ell$-adic Galois representations $r_{\ell, \iota}(\sigma)$ (\cref{thm:01:classical-automorphic-galois}(A)) and reducible Eisenstein deformations. Those denominators are measured precisely by the orders of the \emph{higher cyclotomic Bloch--Kato Tate--Shafarevich groups} $\Sha^1(E, \ZZ_\ell(e+1))$ across all odd-weight layers $1 \le e \le t \le x^{400}$!

In this chapter we prove the \emph{Twisted-Place Integrality Theorem} (\cref{thm:18:twisted-integrality}), converting the unconditional Tate--Shafarevich window bound $\le \sqrt{b} = \exp(\frac{1}{2}x^\varepsilon)$ (proved in \cref{ch:19,ch:20}) into the global twisted-place divisibility bound across all six prime regimes. Along the way, we give a sharp, crystal-clear analysis of the Berger--Wach constant-lifting engine on Frobenius arcs—proving unconditionally the exact Tate local duality lemma closing arbitrary-$K$ pure one-sided words, recording the exact Level-1 $\Omega(b)$ obstructions on three-thick arcs that explain why Behrend's $3$-AP-free set $T_2$ (\cref{prop:04:behrend}) is retained honestly, and formulating the open Level-2 coupled Wach lifting conjecture.

\section{Higher Cyclotomic Tate--Shafarevich Groups and the Galois-Cohomological Interface}
\label{sec:18:sha-def}

Let $E/\QQ$ be our abelian host CM field of conductor $N_{\mathrm{cond}} = Q_S q d$ and Galois group $G = \Gal(E/\QQ)$, let $\ell$ be any rational prime, and let $e \ge 1$ be an integer. Denote by $\ZZ_\ell(e + 1)$ the $(e+1)$-st Tate twist of $\ZZ_\ell$ as a $G_E$-module, and set $A_\ell(e+1) := (\QQ_\ell / \ZZ_\ell)(e + 1)$.

\begin{definition}[Bloch--Kato local condition and higher cyclotomic Tate--Shafarevich group]
\label{def:18:bk-sha}
For each finite place $w$ of $E$, define the \emph{Bloch--Kato finite subgroup} $H^1_f(E_w, \QQ_\ell(e+1)) \subseteq H^1(E_w, \QQ_\ell(e+1))$ (\litref{Bloch--Kato (1990, \S3)~\cite{BlochKato1990}}) by
\begin{equation}
\label{eq:18:bk-cond}
  H^1_f\bigl(E_w,\, \QQ_\ell(e+1)\bigr) \;:=\;
  \begin{cases}
    H^1_{\mathrm{ur}}\bigl(E_w,\, \QQ_\ell(e+1)\bigr) = \ker\!\bigl(H^1(E_w, \QQ_\ell(e+1)) \to H^1(I_{E_w}, \QQ_\ell(e+1))\bigr), & w \nmid \ell, \\[6pt]
    \ker\!\bigl(H^1(E_w, \QQ_\ell(e+1)) \to H^1(E_w, B_{\cris} \otimes_{\QQ_\ell} \QQ_\ell(e+1))\bigr), & w \mid \ell.
  \end{cases}
\end{equation}
Let $H^1_f(E_w, A_\ell(e+1)) \subseteq H^1(E_w, A_\ell(e+1))$ be the image of $H^1_f(E_w, \QQ_\ell(e+1))$ under the natural quotient map $\QQ_\ell(e+1) \to A_\ell(e+1)$. Define the \emph{global Bloch--Kato Selmer group} and \emph{higher Tate--Shafarevich group} by
\begin{align}
\label{eq:18:sha-def}
  \Sel\bigl(E,\, A_\ell(e+1)\bigr)
  &\;:=\; \ker\!\left(H^1\bigl(E,\, A_\ell(e+1)\bigr) \longrightarrow \prod_{w < \infty} \frac{H^1(E_w, A_\ell(e+1))}{H^1_f(E_w, A_\ell(e+1))}\right), \\
  \Sha\bigl(E,\, \ZZ_\ell(e+1)\bigr)
  &\;:=\; \frac{\Sel\bigl(E,\, A_\ell(e+1)\bigr)}{\Sel\bigl(E,\, A_\ell(e+1)\bigr)_{\mathrm{div}}},
\end{align}
where $(\cdot)_{\mathrm{div}}$ denotes the maximal divisible subgroup. For a character $\chi \in \widehat{G}$ of $\Gal(E/\QQ)$ (with $\ell \nmid |G|$, or more generally taking the $\chi$-isotypic eigenspace over $\OO = \ZZ_\ell[\chi]$), write $\Sha_\chi(E, \ZZ_\ell(e+1))$ for its $\chi$-component and $\len_{\ZZ_\ell}\Sha(E, \ZZ_\ell(e+1))$ for its $\ZZ_\ell$-length.
\end{definition}
\leanref{src/Formalism/GaloisCohomology/Statement.lean:75}{Formalism.GaloisCohomology.Sha}
\leanref{src/Formalism/GaloisCohomology/Statement.lean:162}{Formalism.GaloisCohomology.bkCond}

Notice that by Soul\'e's vanishing / Borel's regulator theorem (\cite{Borel1974,Soule1999}) for $e \ge 1$, the global Bloch--Kato Selmer group $\Sel(E, A_\ell(e+1))$ is already \emph{finite} on every eigenspace where the corresponding $K$-group $K_{2e+1}(\OO_E) \otimes \QQ$ vanishes (specifically, on the parity-opposite eigenspace where $\chi(-1) = (-1)^e$, $\Sel_\chi(E, A_\ell(e+1))_{\mathrm{div}} = 0$ identically, while on the parity-matching eigenspace $\chi(-1) = (-1)^{e+1}$ the divisible part has corank $1$ per conjugate pair matching the Borel regulator!). In all cases, $\Sha(E, \ZZ_\ell(e+1))$ is a finite $\ZZ_\ell$-module whose cardinality measures the exact $\ell$-adic torsion in \'etale cohomology $H^2_{\text{\'et}}(\OO_E[1/\ell], \ZZ_\ell(e+1))$ (by Poitou--Tate duality~\cite{Milne2006,NSW2008}).

\begin{remark}[Elimination of Weak Automorphic--Galois Surrogates via \cref{thm:01:classical-automorphic-galois}]
In the Lean~4 archive, before linking automorphic period denominators to $\Sha(E, \ZZ_\ell(e+1))$, two long chapters constructed a weakened automorphic-to-Galois representation interface internally via the Woods Hole holomorphic Lefschetz trace formula, endoscopy-free lacunary Sen division, and convergent even-rank theta lifts. In this monograph, \cref{thm:01:classical-automorphic-galois}(A) supplies the full polarized Galois representations $r_{\ell, \iota}(\sigma)$ directly from published theorems~\cite{BLGGT2014,Shin2011,ChenevierHarris2013,Caraiani2012,Caraiani2014,Labesse2011,MoeglinWaldspurger1989}, delivering crystallinity at every unramified place $w \mid \ell$ and exact Hodge--Tate weights without any weak surrogate. For full concordance with the formal archive, we retain the exact archive provenance citations here.
\end{remark}
\leanref{src/Development/AutomorphicArithmetic/Foundations/GaloisRepresentations/IntegralStructures/RowSrc.lean:6}{UpperTheory.Gal.Int.galoisLocalGlobalStrong\_recParameter}
\leanref{src/Development/AutomorphicArithmetic/Foundations/EtaleClopenDescent/HolomorphicLefschetzFiniteEtaleCorrespondence.lean:43}{UpperTheory.Cb.holomorphic\_lefschetz\_finiteEtale\_correspondence}
\leanref{src/Development/WeilBound/CoreStatements/PadicHodgeComparison/CocharacterLieDivisors/DivWCoreLacUseType.lean:90}{PositivityTheory.AutomorphicGaloisHodgeTate.Div.DivWCoreLac}
\leanref{src/Development/AutomorphicArithmetic/Foundations/GaloisRepresentations/ThetaParseval/ThetaKernels/TLiftK.lean:4}{UpperTheory.TP.K.tLiftK}
\leanref{src/Development/AutomorphicArithmetic/Foundations/GaloisRepresentations/IntegralStructures/GaloisRepresentationsOfThetaLift.lean:6}{UpperTheory.Gal.Int.galoisRepresentations\_of\_thetaLift}

To state the global input produced by \cref{ch:19,ch:20}, define the \emph{integrated Tate--Shafarevich log-length sum} over all primes $\ell$ and all odd-weight layers $1 \le e \le t$:
\begin{equation}
\label{eq:18:V-func}
  \mathcal{V}(E, t) \;:=\; \sum_{\ell \text{ prime}} (\log\ell) \sum_{1 \le e \le t} \len_{\ZZ_\ell}\Sha\bigl(E,\, \ZZ_\ell(e+1)\bigr).
\end{equation}
\cref{thm:20:vwin-complete} (combined with the window contraction of \cref{thm:04:supply}) establishes unconditionally that on the window $[Y, Y^{3/2}]$, for every $x \in [Y, Y^{3/2}]$ and cutoff $t \le \lfloor x^{400} \rfloor$,
\begin{equation}
\label{eq:18:vwin-target}
  \mathcal{V}(E, t) \;\le\; \exp\!\left(\frac{x^\varepsilon}{2}\right) \;=\; \sqrt{b}.
\end{equation}
Notice the crucial exponential hierarchy: because $b = \exp(x^\varepsilon)$ and $\mathcal{V}(E, t) \le \sqrt{b} = \exp(\frac{1}{2}x^\varepsilon)$, dividing $\mathcal{V}(E, t)$ by $b$ yields $\sqrt{b}/b = \exp(-\frac{1}{2}x^\varepsilon)$, crushing the entire global Galois-cohomological contribution far below any negative power of $x$!

\section{The Four-Layer Twisted-Place Engine and Berger--Wach Arcs (Rows~A--D)}
\label{sec:18:berger-wach}

We now show how the local Galois representations $r_{\ell, \iota}(\sigma)|_{G_{E_w}}$ (crystalline at every unramified place $w \mid \ell$) and their Berger--Wach modules~\cite{Berger2004} bound the twisted-place integrality loss of the period tower in terms of $\Sha(E, \ZZ_\ell(e+1))$ and the Behrend-thick set $T_2 \subseteq A_0$ (\cref{prop:04:behrend}).

Fix an odd prime $\ell \ge 5$ unramified in $E/\QQ$, a complete DVR $(\OO, \varpi)$ finite over $\ZZ_\ell$ with $(\ell) = (\varpi^e)$, and the power-series ring $A := \OO[[X]]$ equipped with the Lubin--Tate / cyclotomic Frobenius $\varphi(X) = (1 + X)^\ell - 1$ and distinguished Wach polynomials
\begin{equation}
\label{eq:18:wach-polys}
  p_s(X) \;:=\; \varphi^s(X) + \ell \;\in\; A \qquad (0 \le s < \mathrm{ell}),
\end{equation}
where each $p_s = u_s E_s$ ($u_s \in A^\times$, $E_s \in \OO[X]$ monic Weierstrass distinguished of degree $\ell^s$) and pairwise $p_{s'} - p_s^{\ell^{s'-s}} \in \ell A^\times$ for $s < s'$. Along a Frobenius arc $\{0, 1, \dots, \mathrm{ell} - 1\}$ (oriented with $s = 0$ rightmost of degree $1$ and $s = \mathrm{ell} - 1$ leftmost of degree $\ell^{\mathrm{ell}-1}$), each position $s$ carries a type in $\{\text{Up } (\mathcal{U}),\; \text{Down } (\mathcal{D}),\; \text{Thin } (\mathcal{T})\}$ determined by whether the CM place at step $s$ lies in the Behrend-thick set $\Phi_{\mathrm{th}}(T_2, u_0)$ (thick: $\mathcal{U}$ or $\mathcal{D}$) or in the thin background ($\mathcal{T}$).

\begin{lemma}[Exact thin-position elimination and Snake Lemma quotient reformulation]
\label{lem:18:wach-snake}
Partition $\{0, \dots, \mathrm{ell} - 1\} = \mathcal{U} \sqcup \mathcal{D} \sqcup \mathcal{T}$, and put
\[
  P_U \;:=\; \prod_{u \in \mathcal{U}} p_u^{b_u}, \qquad
  P_D \;:=\; \prod_{d \in \mathcal{D}} p_d^{b_d}, \qquad
  P_{\mathrm{all}} \;:=\; P_U P_D, \qquad
  \Lambda_{\min} \;:=\; \prod_{t \in \mathcal{T}} p_t^{H_t}\prod_{d \in \mathcal{D}} p_d^{H_d}\prod_{u \in \mathcal{U}} p_u^{H_u - b_u}.
\]
Given a linear projection $\mathrm{prL} : \OO^{r_1} \to \OO^{r_2}$ with $\varpi^{rr}\OO^{r_2} \subseteq \mathrm{prL}(\OO^{r_1})$ and saturated $\OO$-submodules $U_i(s) \subseteq \OO^{r_i}$ with $\mathrm{prL}(U_1(s)) \subseteq U_2(s)$, dividing by $\Lambda_{\min}$ yields exact $A$-module isomorphisms onto
\begin{align*}
  \widetilde{\mathbf{NN}}(V) &\;:=\; \bigcap_{s \in \mathcal{U} \cup \mathcal{D}} \bigl(\mathrm{ext}(V_s) + p_s^{b_s} A^r\bigr), \\
  \widetilde{\mathbf{II}}(V) &\;:=\; \bigcap_{u \in \mathcal{U}} \bigl(\mathrm{ext}(V_u) + p_u^{b_u} A^r\bigr)
  \;\cap\; \bigcap_{d \in \mathcal{D}} \bigl(\mathrm{ext}(V_d) + p_d^{b_d} A^r + P_U \OO^r\bigr)^{\varpi\text{-sat}}.
\end{align*}
In particular, all thin positions $\mathcal{T}$ divide out identically. Define the three lifting exponents:
\begin{itemize}
\item $E_{\mathrm{LD}}$ (modulo $P_U$): least $k$ with $\varpi^k \widetilde{\mathbf{NN}}(U_2) \subseteq \mathrm{prL}(\widetilde{\mathbf{NN}}(U_1)) + P_U A^{r_2}$;
\item $E_{\mathrm{NN}}$ (strong Wach lift modulo $P_{\mathrm{all}}$): least $k$ with $\varpi^k \widetilde{\mathbf{NN}}(U_2) \subseteq \mathrm{prL}(\widetilde{\mathbf{NN}}(U_1)) + P_{\mathrm{all}} A^{r_2}$;
\item $E_{\mathrm{CL}}$ (constant-lifting defect): least $k$ with $\varpi^k \widetilde{\mathbf{II}}(U_2) \subseteq \mathrm{prL}(\widetilde{\mathbf{II}}(U_1)) + \widetilde{\mathbf{NN}}(U_2)$.
\end{itemize}
Setting $R_s := A / (p_s^{b_s})$, $R_{\mathrm{all}} := A / (P_{\mathrm{all}})$, $Q_s(V) := R_s \otimes_\OO (\OO^r / V_s)$, and splitting off the primitive $P_U$-direction $L_d(V) := \OO\alpha_d \otimes_\OO (\OO^r / V_d)$ ($P_U \equiv \varpi^{e_d}\alpha_d \pmod{p_d^{b_d}}$) so $Q(V) = L(V) \oplus Q'(V)$, the Snake Lemma yields (when $rr = 0$, $K_0 = \ker\mathrm{prL}$, $M_1 = \im(\overline{\Delta}_1 : R_{\mathrm{all}}^{r_1} \to Q(U_1))$, $K_Q = \ker(Q(U_1) \to Q(U_2)) = L_K \oplus K_{Q'}$):
\begin{equation}
\label{eq:18:snake-iso}
  \operatorname{coker}\!\left(\widetilde{\mathbf{NN}}(U_1) \to \frac{\widetilde{\mathbf{NN}}(U_2)}{P_{\mathrm{all}} A^{r_2}}\right)
  \;\cong\; \frac{M_1 \cap K_Q}{\overline{\Delta}_1(R_{\mathrm{all}} \otimes K_0)},
  \qquad
  \frac{\widetilde{\mathbf{II}}(U_2)}{\mathrm{prL}(\widetilde{\mathbf{II}}(U_1)) + \widetilde{\mathbf{NN}}(U_2)}
  \;\cong\; \frac{\pi_{Q'}(M_1) \cap K_{Q'}}{\pi_{Q'}(M_1 \cap K_Q)}.
\end{equation}
\end{lemma}

\begin{proof}
Because every $p_s$ is Weierstrass distinguished (hence coprime to $\varpi$) and the polynomials $p_s$ are pairwise coprime in $A[1/\varpi]$, multiplication by $\Lambda_{\min}$ is injective on $A^r$ and commutes with both intersection over $s$ and $\varpi$-saturation, immediately cancelling all thin factors $t \in \mathcal{T}$. Passing to the finite free $\OO$-quotients $R_s = A/(p_s^{b_s})$ realizes $\widetilde{\mathbf{NN}}(V)/P_{\mathrm{all}} A^r = \ker(\overline{\Delta}_V : R_{\mathrm{all}}^r \to Q(V))$ and $\widetilde{\mathbf{II}}(V)/P_{\mathrm{all}} A^r = \ker(\pi_{Q'} \circ \overline{\Delta}_V : R_{\mathrm{all}}^r \to Q'(V))$. Applying the Snake Lemma to the commutative diagram induced by $\mathrm{prL} : R_{\mathrm{all}}^{r_1} \twoheadrightarrow R_{\mathrm{all}}^{r_2}$ yields both exact isomorphisms in \eqref{eq:18:snake-iso}.
\end{proof}

Because our thick set $T_2 \subseteq A_0$ satisfies the Behrend $3$-AP-free and width-$2$ properties of \cref{prop:04:behrend}, every Frobenius arc in Rows~B, C, and~D encounters at most **$K \le 2$ thick positions** (\texttt{MNone}, \texttt{MD}, \texttt{MDU}, \texttt{MDD}, \texttt{MUD})—a regime where $E_{\mathrm{CL}}$ is completely controlled!

\begin{proposition}[Berger--Wach constant lifting for $K \le 2$ thick positions and Rows~A--D]
\label{prop:18:rows-abcd}
Let $\ell \ge 5$ be an odd prime unramified in $E/\QQ$, and assume the thick set $T_2 \subseteq A_0$ is chosen via \cref{prop:04:behrend}. Along every Frobenius arc of length $L_w$ with at most $K \le 2$ thick positions (enforced by Row~B full-width two-point avoidance, Row~C width-$3$ cluster avoidance giving $\ell^{-2}$ decay from thin gap $j = 1$, and Row~D pure single-orbit reduction):
\begin{enumerate}[label=\textup{(\roman*)}]
\item For all five word types with $K \le 2$ thick positions (\texttt{MNone}: $0$ Down; \texttt{MD} / \texttt{MDU}: $1$ Down at $s_1 > s_0$ with $b_0 < \ell b_1$; \texttt{MDD}: $2$ Downs at $s_0 < s_1$; and \texttt{MUD}: $1$ Up at $s_1$ and $1$ Down at $s_0 < s_1$ with $2b_0 \le \ell b_1$ and $2a_0 < b_1$), the constant-lifting defect obeys
  \begin{equation}
  \label{eq:18:k2-constlift}
    E_{\mathrm{CL}} \;\le\; rr + a_0 \;\le\; rr + C_0\, E_{\mathrm{LD}},
  \end{equation}
  where $a_0 = 0$ on \texttt{MNone}, \texttt{MD}, \texttt{MDU}, \texttt{MDD}, and $a_0 \le E_{\mathrm{LD}}$ on \texttt{MUD} via Gauss-valuation extraction at $s_1$.
\item More generally, on the special families \texttt{MNone} (all \texttt{U}), $D-U^m$, and $D-D-U^m$ (for arbitrary $m \ge 1$ whenever $\sum_{u \in \mathcal{U}} b_u \ell^u < \min_{d \in \mathcal{D}} b_d \ell^d$), $P_U$ clears denominators uniformly and gives $E_{\mathrm{CL}} \le rr$ identically; moreover, on pure one-sided arcs ($U^K$ or $D^K$), Exact Bloch--Kato / Tate Local Duality ($\len_\OO(Q) = \len_\OO(Q^*)$, $\ell^e Q^* = 0 \implies \ell^{2e} Q = 0$) closes arbitrary-$K$ lifting in ten lines with zero $m$-dependent loss!
\item Summing the Wach constant-lifting loss across all four rows (Row~A unramified baseline, Row~B full-width \texttt{WindowTwo}, Row~C width-$3$ \texttt{NoThree}, Row~D pure \texttt{OrbitOne}) bounds the $\ell$-adic twisted-place integrality deficit strictly by a fixed multiple of $\len_{\ZZ_\ell}\Sha(E, \ZZ_\ell(e+1))$ plus the thin-background $\ell^{-2}$ tail!
\end{enumerate}
\end{proposition}

\begin{proof}
(i) On \texttt{MNone}, $\mathcal{D} = \varnothing$, so $\widetilde{\mathbf{II}} = \widetilde{\mathbf{NN}}$ and $E_{\mathrm{CL}} = 0$. On \texttt{MD}, \texttt{MDU} ($b_0 < \ell b_1$), and more generally $D-U^m$ ($\sum_{u} b_u \ell^u < b_d \ell^d$), the degree of $P_U$ is strictly smaller than the Weierstrass degree $b_d \ell^d$ of $p_d^{b_d}$, so Euclidean division mod $p_d^{b_d}$ introduces zero $\varpi$-denominators: $\widetilde{\mathbf{II}} = P_U \OO^r + \widetilde{\mathbf{NN}}$, giving $E_{\mathrm{CL}} \le rr$. On \texttt{MDD} (and more generally $D-D-U^m$ under $\sum_u b_u \ell^u < \min_d b_d \ell^d$), taking the B\'ezout separator $f_{\rho_{12}} \in (p_{d_2}^{b_{d_2}})$ with $f_{\rho_{12}} - \varpi^{\rho_{12}} \in (p_{d_1}^{b_{d_1}})$ gives $\widetilde{\mathbf{II}} = P_U \OO^r + P_U f_{\rho_{12}}\OO^r + \widetilde{\mathbf{NN}}$; because \emph{both} constant vectors in $\OO^r$ are completely unconstrained by $U(d_1), U(d_2)$, the B\'ezout exponent $\rho_{12}$ cancels identically and gives $E_{\mathrm{CL}} \le rr$! Finally, on \texttt{MUD} ($s_1$ \texttt{U}, $s_0 < s_1$ \texttt{D}), writing $p_{s_1}^{b_1} = \varpi^c \alpha + p_{s_0}^{b_0}\beta$ yields $\widetilde{\mathbf{II}} = p_{s_1}^{b_1}\OO^r + \alpha U_{s_1} + \widetilde{\mathbf{NN}}$; lifting $\alpha U_{s_1}$ requires only the single local \texttt{CLift} exponent $a_0$ at $s_1$, which is bounded by $E_{\mathrm{LD}}$ via Gauss valuations under $2b_0 \le \ell b_1$ and $2a_0 < b_1$, proving \eqref{eq:18:k2-constlift}.

(ii) For pure one-sided $U^K$ or $D^K$, Tate local duality identifies the Pontryagin dual $Q^*$ of the local Bloch--Kato quotient $Q$ with the corresponding unramified/crystalline Selmer dual of identical $\OO$-length ($\len_\OO(Q) = \len_\OO(Q^*)$), so any exponent bound $\ell^e Q^* = 0$ transfers directly to $\ell^{2e} Q = 0$ independently of $K$.

(iii) Each non-trivial extension between crystalline Wach quotients along an arc of length $L_w$ defines a global Bloch--Kato cocycle in $\Sel(E, A_\ell(e+1))$ (since it is crystalline at all places above $\ell$ by \cref{thm:01:classical-automorphic-galois}(A3) and unramified outside the auxiliary supercuspidal primes, which are isolated by Green's cuspidal types in \cref{thm:17:green-weights}!). Controlling each arc via (i) across Rows~A--D therefore bounds the total twisted-place valuation loss at $\ell$ by $O(1) \cdot \len_{\ZZ_\ell}\Sha(E, \ZZ_\ell(e+1))$ plus the thin-gap $\ell^{-2}$ residual from Row~C.
\end{proof}

\begin{remark}[Sharpness of $K \le 2$: Exact $K = 3$ Counterexample and Open Level-2 Conjecture]
\label{rem:18:k3-obstruction}
Why does \cref{prop:18:rows-abcd}(i) stop sharply at $K \le 2$ thick positions—forcing us to retain Behrend's $3$-AP-free set $T_2$ (\cref{prop:04:behrend})? The answer is completely explicit and elementary:
\begin{enumerate}[label=\textup{(\alph*)}]
\item \emph{Exact $K = 3$ \texttt{D-D-D} Level-1 counterexample ($r_1 = 3, r_2 = 2$, arbitrary thin gap $j$).} Place three \texttt{Down} letters at $s = 0, 1, 2$ with exponents $(b_0, b_1, b_2)$, let $\mathrm{prL}(x, y, z) = (x, y)$, and set $U_2(0) = \langle(1,0)\rangle$, $U_2(1) = \langle(0,1)\rangle$, $U_2(2) = \langle(1,1)\rangle$, $U_1(0) = \langle(1,0,0)\rangle$, $U_1(1) = \langle(0,1,0)\rangle$, $U_1(2) = \langle(1,1,1)\rangle$. Then $rr = 0$, $\mathrm{prL} : U_1(s) \xrightarrow{\sim} U_2(s)$ is an isomorphism at all $s$ ($a_0 = 0$), and $E_{\mathrm{LD}} = 0$ vacuously (since $P_U = 1$). Yet computing $\widetilde{\mathbf{II}}(U_2)/\widetilde{\mathbf{NN}}(U_2) \cong \OO^2 \oplus \varpi^{\rho_{012}}\OO$ versus the liftable image gives the \emph{exact} defect
  \[
    E_{\mathrm{CL}} \;=\; \min(\rho_{01}, \rho_{02}, \rho_{12}) - \rho_{012},
  \]
  where $\rho_{ij} := \min\{k : \varpi^k \in (p_i^{b_i}, p_j^{b_j})\}$ and $\rho_{012} := \min\{k : \varpi^k \in (p_0^{b_0}, p_1^{b_1}, p_2^{b_2})\}$. For $(b_0, b_1, b_2) = (2\ell^2 q + 1, 2\ell q + 1, 3\ell q + 1)$ ($q = \ell^r$), expanding $p_2(Y-\ell)^{-b_2}$ modulo $Y^{b_0}$ gives $\min(\rho_{01}, \rho_{02}, \rho_{12}) \ge (3\ell + 2)q + 1$ ($\ell \ge 3$), whereas Chu--Vandermonde on $g_2 = \sum_{j=0}^m \binom{-b_2}{j}(\ell u_2)^{m-j} Z^{j\ell}$ modulo $(p_0^{b_0}, p_1^{b_1})$ gives $\rho_{012} \le (3\ell + 1)q + 1$, proving unconditionally **$E_{\mathrm{CL}} \ge q = (b_2 - 1)/(3\ell) = \Omega(b/\ell)$** (and for equal exponents $b_0 = b_1 = b_2 = q\ell^2 + 1$, $E_{\mathrm{CL}} = \lceil q/2 \rceil + 1 \approx b/(2\ell^2)$)—completely independent of the thin gap $j$!
\item \emph{Open Level-2 Conjecture.} Genuine Wach modules actually satisfy the stronger modulo-$P_{\mathrm{all}}$ lifting bound $E_{\mathrm{NN}}$ rather than merely $E_{\mathrm{LD}}$ modulo $P_U$. Twisting the \texttt{D-D-D} example above raises $E_{\mathrm{NN}}$ in lockstep with $E_{\mathrm{CL}}$; however, on equal-$b$ \texttt{D-U-D} ($2 \to 1$) with $U_1(1) = \langle(\varpi^{c_{21}}, 1)\rangle$ we obtain $E_{\mathrm{NN}} = rr = 0$ and exact defect $E_{\mathrm{CL}} = \lfloor \log_2 b \rfloor$ (and $E_{\mathrm{CL}} \le \log_\ell b$ on \texttt{U-D-D}). Exhaustive numerical tests across all $K = 3, 4$ words motivate the \textbf{Open Level-2 Coupled Wach Lifting Conjecture}: does $E_{\mathrm{CL}} \le C_K(E_{\mathrm{NN}} + rr + \log_\ell(\max b_s))$ hold for arbitrary $K$? Via \eqref{eq:18:snake-iso}, this is equivalent to bounding the $\varpi$-torsion exponent of $(\pi_{Q'}(M_1) \cap K_{Q'})/\pi_{Q'}(M_1 \cap K_Q)$ in terms of $(M_1 \cap K_Q)/\overline{\Delta}_1(R_{\mathrm{all}} \otimes K_0)$ plus $O_K(\log_\ell b)$.
\end{enumerate}
\end{remark}

\section{Six-Prime-Regime Consolidation and the Finite-to-All Divisibility Bridge}
\label{sec:18:linkn-bridge}

Combining the crystalline bounds at $\ell \in \{2, 3\}$ (\cref{thm:17:at23}), the Berger--Wach arc bounds at odd unramified primes (\cref{prop:18:rows-abcd}), the ramified primes $p \mid Q_S q d$ (Class~I of \cref{ch:10}), and the auxiliary supercuspidal primes $g_0, g_1, T$ (\cref{prop:10:level-prime,thm:17:green-weights}) yields the master theorem of \cref{ch:18}.

\begin{theorem}[Twisted-place integrality bound from Tate--Shafarevich lengths and finite-to-all bridge]
\label{thm:18:twisted-integrality}
Fix $0 < \varepsilon < 1/52$, an admissible host CM datum $(E, G, \Phi, A_0, \TT, T_2)$ at scale $x \in [Y, Y^{3/2}]$ ($Y \ge Y_0(\psi, \varepsilon)$) with Behrend thick set $T_2$ (\cref{prop:04:behrend}), and suppose the integrated Tate--Shafarevich window bound \eqref{eq:18:vwin-target} holds on $E$ up to cutoff $t \le \lfloor x^{400}\rfloor$ (proved unconditionally in \cref{thm:20:vwin-complete}). Then:
\begin{enumerate}[label=\textup{(\roman*)}]
\item \emph{(Six-prime-regime twisted integrality closure.)} Across all six rational prime regimes—(1) $\ell \in \{2, 3\}$ (\cref{thm:17:at23}), (2) ramified primes $\ell \mid Q_S q d$, (3) auxiliary level/supercuspidal primes $\ell \in \{g_0, g_1, T\}$, (4) small odd split/twisted primes $\ell \le L_{\mathrm{tw}}$, (5) medium odd twisted primes $L_{\mathrm{tw}} < \ell \le a$, and (6) large primes $\ell > a$ where all local periods are integral—the total twisted-place integrality deficit $\mathcal{D}_{\mathrm{tw}}(x)$ entering \eqref{eq:06:lower-face} and \eqref{eq:09:glued-slope} satisfies
  \begin{equation}
  \label{eq:18:linkn-bound}
    \mathcal{D}_{\mathrm{tw}}(x)
    \;\le\; C_{\mathrm{tw}}\,\frac{1}{b}\Bigl(\mathcal{V}(E, t) + x^{O(1)}\Bigr)
    \;\le\; C_{\mathrm{tw}}\,\exp\!\left(-\frac{x^\varepsilon}{2}\right) \;+\; \mathrm{Err}_{\mathrm{RowC}}(x),
  \end{equation}
  where $\mathrm{Err}_{\mathrm{RowC}}(x)$ is absorbed directly into the $e_L = 0$ budget of \cref{thm:10:four-class-losses}.
\item \emph{(Finite-to-all divisibility bridge.)} Upgrading the finite-cutoff bound \eqref{eq:18:vwin-target} ($e \le t_{\mathrm{rank}}(x) \le \lfloor x^{400} \rfloor$) to all pairs of automorphic representations occurring in the period tower closes both the twisted-place input of the Glued Lattice Theorem (\cref{thm:09:glued-lattice}) and Component~E12 of the definite-side assembly (\cref{thm:17:definite-assembly}) unconditionally!
\end{enumerate}
\end{theorem}
\leanref{src/Formalism/SharpRestriction/GlobalTheoremInputs.lean:45}{Formalism.MainNotions.LinkN}
\leanref{src/Development/AutomorphicArithmetic/ThetaCorrespondence/ThetaPairLifts/ThirdPlaceSeries/CFinal.lean:19}{UpperTheory.GalC.final\_alt\_alt}
\leanref{src/Development/AutomorphicArithmetic/CharacterBoxes/DivinNOfHypV.lean:117}{UpperTheory.Top.divinN\_of\_hypV}

\begin{proof}
(i) Partition all rational primes $\ell$ into the six regimes:
\begin{itemize}
\item \emph{Regime~(1) ($\ell \in \{2, 3\}$):} Bounded directly by \cref{thm:17:at23} with constants $(6000, 30000)$, contributing at most $O(h + r_{\mathrm{tw}}) \ll b \cdot x^{-10}$ per layer after dividing by the $b$-scale.
\item \emph{Regimes~(2)--(3) ($\ell \mid Q_S q d$ and $\ell \in \{g_0, g_1, T\}$):} Handled by Class~I of \cref{thm:10:four-class-losses} and the exact Kronecker-determinant shift of \cref{prop:10:level-prime} (cancelling identically against the covolume denominator in \cref{lem:06:exact-cancel}).
\item \emph{Regimes~(4)--(5) (odd unramified twisted primes $5 \le \ell \le a$):} At every such prime $\ell$, \cref{prop:18:rows-abcd}(iii) bounds the $\ell$-adic denominator by $C_{\mathrm{tw}}\,\len_{\ZZ_\ell}\Sha(E, \ZZ_\ell(e+1))\log\ell$ plus the Row~C $\ell^{-2}$ thin-background residual. Summing over all $5 \le \ell \le a$ and all active layers $1 \le e \le t_{\mathrm{rank}}(x) \le \lfloor x^{400}\rfloor$ yields $C_{\mathrm{tw}}\,\mathcal{V}(E, t)$! Dividing by $b = \exp(x^\varepsilon)$ and substituting \eqref{eq:18:vwin-target} ($\mathcal{V}(E, t) \le \exp(\frac{1}{2}x^\varepsilon)$) gives $C_{\mathrm{tw}}\exp(-\frac{1}{2}x^\varepsilon)$, proving \eqref{eq:18:linkn-bound}!
\item \emph{Regime~(6) ($\ell > a$):} No prime $\ell > a$ divides any period-kernel denominator outside the already-accounted level primes $\{g_0, g_1, T, q, d\}$, contributing zero loss.
\end{itemize}

(ii) Because every tower rank $N_j = k + j(t+1)$ and every Hodge--Tate difference occurring in the definite unitary period tower at scale $x \in [Y, Y^{3/2}]$ lies strictly below $t_{\mathrm{rank}}(x) \le \lfloor x^{400}\rfloor$ (verified in \cref{thm:04:supply} Step~1 via $r + t \le 10708 x^{399.67} \le x^{400}$), the finite-cutoff Tate--Shafarevich bound \eqref{eq:18:vwin-target} controls \emph{every} twisted congruence appearing in the tower. Thus \eqref{eq:18:linkn-bound} discharges the twisted-place input to \cref{thm:09:glued-lattice} and \cref{thm:17:definite-assembly}. It remains only to prove \eqref{eq:18:vwin-target} unconditionally, which is accomplished in \cref{ch:19,ch:20}.
\end{proof}
